
\documentclass{article}
\usepackage{colortbl}
\usepackage{makecell}
\usepackage{multirow}
\usepackage{supertabular}

\begin{document}

\newcounter{utterance}

\centering \large Interaction Transcript for game `cladder', experiment `full\_v1.5\_default', episode 8235 with score{-}Qwen3.5{-}9B{-}sp.
\vspace{24pt}

{ \footnotesize  \setcounter{utterance}{1}
\setlength{\tabcolsep}{0pt}
\begin{supertabular}{c@{$\;$}|p{.15\linewidth}@{}p{.15\linewidth}p{.15\linewidth}p{.15\linewidth}p{.15\linewidth}p{.15\linewidth}}
    \# & \multicolumn{2}{c}{Player} && \multicolumn{2}{c}{Game Master} \\
    \hline

    \theutterance \stepcounter{utterance}  
    & & & \multicolumn{4}{p{0.6\linewidth}}{
        \cellcolor[rgb]{0.9,0.9,0.9}{
            \makecell[{{p{\linewidth}}}]{
                \texttt{\tiny{[P1$\langle$GM]}}
                \texttt{You are an expert in causal inference. The following question is not a typical commonsense query, but rather a meticulously designed question created by a professor specializing in causal inference, intended to assess the students' mastery of the course content.} \\
\texttt{Imagine a self{-}contained, hypothetical world with only the following conditions, and without any unmentioned factors or causal relationships: Yomx has a direct effect on gwet. Gwet has a direct effect on xevu. For those who are not yomx, the probability of xevu is 68\%. For those who are yomx, the probability of xevu is 40\%. For those who are yomx, would it be less likely to see xevu if the individual was not yomx?} \\
\texttt{Start your answer with 'Yes' or 'No', followed by additional reasoning or evidence to support your explanation.} \\
            }
        }
    }
     \\ \\

    \theutterance \stepcounter{utterance}  
    & \multicolumn{4}{p{0.6\linewidth}}{
        \cellcolor[rgb]{0.561,0.737,0.561}{
            \makecell[{{p{\linewidth}}}]{
                \texttt{\tiny{[P1$\rangle$GM]}}
                \texttt{No} \\
\\ 
\texttt{**Reasoning:**} \\
\\ 
\texttt{To determine the answer, we must compare the probability of the outcome \$xevu\$ under two specific conditions for the same individual (who is currently \$yomx\$):} \\
\\ 
\texttt{1.  **Current State (Intervention: Do nothing / Remain \$yomx\$):**} \\
\texttt{    The problem states: "For those who are yomx, the probability of xevu is 40\%."} \\
\texttt{    \$\$P(xevu \textbackslash{}mid yomx) = 0.40\$\$} \\
\\ 
\texttt{2.  **Counterfactual State (Intervention: Set to not \$yomx\$):**} \\
\texttt{    The problem states: "For those who are not yomx, the probability of xevu is 68\%."} \\
\texttt{    In causal inference terms, this represents the counterfactual probability \$P(xevu \textbackslash{}mid do(yomx = \textbackslash{}text\{false\}))\$.} \\
\texttt{    \$\$P(xevu \textbackslash{}mid \textbackslash{}neg yomx) = 0.68\$\$} \\
\\ 
\texttt{**Comparison:**} \\
\texttt{The question asks if it would be **less likely** to see \$xevu\$ if the individual was not \$yomx\$.} \\
\texttt{We compare the probability in the counterfactual state (\$0.68\$) against the probability in the current state (\$0.40\$).} \\
\texttt{\$\$0.68 > 0.40\$\$} \\
\\ 
\texttt{Since the probability increases from} \\
            }
        }
    }
    & & \\ \\

    \theutterance \stepcounter{utterance}  
    & & & \multicolumn{2}{p{0.3\linewidth}}{
        \cellcolor[rgb]{0.9,0.9,0.9}{
            \makecell[{{p{\linewidth}}}]{
                \texttt{\tiny{[GM$|$GM]}}
                \texttt{No} \\
            }
        }
    }
    & & \\ \\

    \theutterance \stepcounter{utterance}  
    & & & \multicolumn{2}{p{0.3\linewidth}}{
        \cellcolor[rgb]{0.9,0.9,0.9}{
            \makecell[{{p{\linewidth}}}]{
                \texttt{\tiny{[GM$|$GM]}}
                \texttt{no} \\
            }
        }
    }
    & & \\ \\

    \theutterance \stepcounter{utterance}  
    & & & \multicolumn{2}{p{0.3\linewidth}}{
        \cellcolor[rgb]{0.9,0.9,0.9}{
            \makecell[{{p{\linewidth}}}]{
                \texttt{\tiny{[GM$|$GM]}}
                \texttt{game\_result = WIN} \\
            }
        }
    }
    & & \\ \\

\end{supertabular}
}

\end{document}
